% !TeX program = pdflatex
\documentclass[10pt]{beamer}

%========================
% Tema y estilo
%========================
\usetheme{Madrid}
\setbeamertemplate{navigation symbols}{}
\setbeamertemplate{footline}[frame number]

%========================
% Idioma y tipografía
%========================
\usepackage[spanish,es-nodecimaldot]{babel}
\usepackage[utf8]{inputenc}
\usepackage[T1]{fontenc}
\usepackage{lmodern}

%========================
% Matemáticas
%========================
\usepackage{amsmath,amssymb,mathtools}

%========================
% Metadatos
%========================
\title[Salud en México]{Parte II (México): Salud, informalidad y espacio fiscal}
\subtitle{Seguridad social contributiva, paquete diferenciado y seguros privados}
\author{Curso de Economía Poblacional}
\institute{Tecnológico de Monterrey}
\date{\today}

\begin{document}

%------------------------------------------------
\begin{frame}
\titlepage
\end{frame}

%------------------------------------------------
\begin{frame}{Ruta y tesis de la Parte II}
\textbf{Tesis:} México enfrenta un triángulo de tensión:
\begin{enumerate}
  \item Seguridad social contributiva (IMSS/ISSSTE)
  \item Informalidad con paquete diferenciado (no contributivo)
  \item Mercado privado (seguros + CAF)
\end{enumerate}

\bigskip
Bajo una restricción central:
\[
\textbf{Espacio fiscal limitado}
\]

\bigskip
\emph{La política viable no es “gastar más” a secas: es reorganizar riesgo e incentivos con un sendero fiscal creíble.}
\end{frame}

%================================================
% ACTO I: Restricción y diagnóstico
%================================================

\begin{frame}{Acto I: La restricción fiscal manda}
La política de salud compite con otras presiones del Estado:
\[
G = G_{\text{pensiones}} + G_{\text{salud}} + G_{\text{otros}}
\]

Restricción presupuestaria (estilizada):
\[
G \le R + \Delta D
\]
donde \(R\) son ingresos públicos y \(\Delta D\) es variación de deuda compatible con sostenibilidad.

\bigskip
\begin{itemize}
  \item Si \(G_{\text{pensiones}}\uparrow\), el margen para \(G_{\text{salud}}\) se estrecha.
  \item El problema de salud debe plantearse como \textbf{reforma bajo restricción fiscal}.
\end{itemize}
\end{frame}

\begin{frame}{El problema estructural en una frase}
\centering
\Large
México promete protección amplia\\
con financiamiento fragmentado\\
y un mercado laboral altamente informal.
\end{frame}

\begin{frame}{Variables de equilibrio del sistema: \(\phi\) y \(\beta\)}
Usaremos dos variables resumen para conectar diseño e incentivos:

\bigskip
\begin{itemize}
  \item \(\phi\): proporción de trabajadores \textbf{formales} (tamaño del pool contributivo).
  \item \(\beta\): exposición de los hogares a \textbf{gasto de bolsillo} (OOP).
\end{itemize}

\bigskip
Identidades mínimas:
\[
G = G_{\text{público}} + OOP
\qquad\text{y}\qquad
OOP = \beta M
\]
donde \(M\) es el costo del evento de salud relevante.

\bigskip
\emph{La política de salud busca subir \(\phi\) (pooling) y bajar \(\beta\) (protección financiera), sin romper la restricción fiscal.}
\end{frame}

%================================================
% ACTO II: Seguridad social contributiva
%================================================

\begin{frame}{Acto II: Seguridad social contributiva (Bismarck real)}
Financiamiento contributivo (estilizado):
\[
G_{\text{contrib}} = \phi \,\tau\, \bar w
\]
donde \(\tau\) es la tasa contributiva y \(\bar w\) el salario promedio formal.

\bigskip
\begin{itemize}
  \item \(\phi\) bajo \(\Rightarrow\) pool pequeño \(\Rightarrow\) presión financiera.
  \item \(\tau\) alto puede aumentar ingresos hoy, pero puede reducir \(\phi\) mañana (margen/informalidad).
\end{itemize}

\bigskip
\emph{En México, Bismarck no falla por teoría: falla por base contributiva estrecha.}
\end{frame}

\begin{frame}{Incentivos: por qué la informalidad importa para salud}
El individuo compara utilidad esperada entre formal e informal.

Formal (con cobertura \(\alpha_F\), contribución \(\tau\) y costo fijo \(\kappa\)):
\[
EU_F=(1-p)u(y(1-\tau)-\kappa)+p\,u(y(1-\tau)-\kappa-(1-\alpha_F)M)
\]

Informal (con cobertura pública limitada \(\alpha_I\) y gasto directo \(\beta\)):
\[
EU_I=(1-p)u(y)+p\,u(y-\beta M)
\]

\bigskip
\textbf{Regla:} elige formal si \(EU_F \ge EU_I\).

\bigskip
\emph{La estructura del paquete para informales afecta \(\phi\).}
\end{frame}

%================================================
% ACTO III: Paquete diferenciado para informales
%================================================

\begin{frame}{Acto III: Paquete diferenciado (el dilema inevitable)}
En la práctica hay dos paquetes:

\begin{itemize}
  \item \textbf{Formal:} paquete relativamente amplio (\(\alpha_F\) alto) financiado por contribuciones.
  \item \textbf{Informal:} paquete más limitado (\(\alpha_I\) menor) financiado por impuestos/transferencias.
\end{itemize}

\bigskip
La tensión:
\[
\alpha_F - \alpha_I \quad \text{define simultáneamente} \quad
\text{equidad y formalización}
\]

\bigskip
\begin{itemize}
  \item Si \(\alpha_I \uparrow\) mucho, puede caer el incentivo a formalizar (\(\phi\downarrow\)).
  \item Si \(\alpha_I\) es muy bajo, sube \(\beta\) y aumenta pobreza/inequidad.
\end{itemize}
\end{frame}

\begin{frame}{Una política realista: paquete básico explícito}
Bajo restricción fiscal, una política viable es \textbf{hacer explícito} un paquete:

\begin{itemize}
  \item \textbf{Paquete básico universal} (prioridades costo-efectivas, prevención, crónicas esenciales).
  \item \textbf{Cobertura ampliada contributiva} para el formal (más servicios/tiempos/tecnologías).
\end{itemize}

\bigskip
Esto se puede expresar como:
\[
\alpha_I = \alpha_0 \quad \text{(básico universal)}
\qquad \text{y}\qquad
\alpha_F = \alpha_0 + \Delta\alpha
\]

\bigskip
\emph{La clave es transparentar el contrato social y alinear incentivos.}
\end{frame}

%================================================
% ACTO IV: Mercado privado y seguros
%================================================

\begin{frame}{Acto IV: El privado como válvula de escape (y riesgo de segmentación)}
Cuando el público no cubre bien, los hogares migran a privado por acceso/tiempos.

Opciones privadas típicas:
\begin{itemize}
  \item Seguro privado tradicional (alto costo, cobertura alta).
  \item Seguro de bajo costo / planes delgados (prima baja, cobertura limitada).
  \item CAF + pago directo “organizado”.
\end{itemize}

\bigskip
\emph{El privado puede bajar \(\beta\) en eventos pequeños, pero puede segmentar el pooling del sistema.}
\end{frame}

\begin{frame}{Modelo: tercera opción (seguro privado)}
Ampliamos el menú del agente:
\[
\max\{EU_F,\,EU_I,\,EU_P\}
\]

Seguro privado:
\[
EU_P=(1-p)u(y-\pi_p)+p\,u(y-\pi_p-(1-\alpha_p)M)
\]

\bigskip
\begin{itemize}
  \item Tradicional: \(\pi_p\) alta, \(\alpha_p\) alta.
  \item Bajo costo: \(\pi_p\) baja, \(\alpha_p\) baja (copagos/límites).
\end{itemize}

\bigskip
\emph{La política debe decidir: ¿complemento regulado o sustituto segmentador?}
\end{frame}

\begin{frame}{Selección adversa y sistema dual}
Si el privado capta población de menor riesgo o mayor ingreso:
\[
p_{\text{pool público}} \uparrow
\Rightarrow
\text{costo esperado público} \uparrow
\]

Y si el público se vuelve “residual”:
\[
\beta \uparrow \quad \text{(para vulnerables)} \qquad
\phi \downarrow \quad \text{(por menor atractivo formal)}
\]

\bigskip
\emph{Sin coordinación, el sistema converge a dualidad: público de alto riesgo + privado de bajo riesgo.}
\end{frame}

%================================================
% ACTO V: Policy bajo espacio fiscal
%================================================

\begin{frame}{Acto V: Policy bajo espacio fiscal (qué SÍ es creíble)}
Bajo poco espacio fiscal, la policy creíble tiene tres rasgos:

\bigskip
\textbf{(1) Priorizar reducción de \(\beta\) rápidamente}
\begin{itemize}
  \item medicamentos, abasto, atención primaria efectiva
  \item protección ante eventos catastróficos (topes OOP)
\end{itemize}

\bigskip
\textbf{(2) Reasignación + eficiencia antes que “nivel”}
\begin{itemize}
  \item compra estratégica, integración de padrones, evitar duplicidades
\end{itemize}

\bigskip
\textbf{(3) Sendero fiscal gradual y defendible}
\[
\Delta G_{\text{salud}} \text{ pequeño pero sostenido}
\]
\end{frame}

\begin{frame}{El “mapa” de política: tres instrumentos}
Instrumentos (estilizados) y objetivos:

\bigskip
\begin{tabular}{p{0.34\textwidth} p{0.55\textwidth}}
\textbf{Instrumento} & \textbf{Objetivo principal} \\
\hline
Paquete básico universal (\(\alpha_0\)) & bajar \(\beta\), legitimidad social \\
Diferencial contributivo (\(\Delta\alpha\)) & sostener \(\phi\), evitar desaliento formal \\
Regulación/integ. privado & evitar segmentación, cobertura complementaria \\
\end{tabular}

\bigskip
\emph{No es “todo para todos”: es un diseño explícito de contrato social bajo restricción fiscal.}
\end{frame}

\begin{frame}{Integración final: el equilibrio \(\phi\)--\(\beta\) bajo restricción}
Base pública:
\[
G_{\text{público}} = \phi\tau \bar w + T
\]
y exposición financiera:
\[
OOP=\beta M
\]

\bigskip
\textbf{Meta de política (bajo restricción fiscal):}
\[
\beta \downarrow \quad \text{sin inducir} \quad \phi \downarrow
\]

\bigskip
\textbf{Lectura:}
\begin{itemize}
  \item bajar \(\beta\) sin destruir \(\phi\) requiere paquete básico + diferencial contributivo claro
  \item el privado debe ser complemento regulado, no sustituto que vacía el pool
\end{itemize}
\end{frame}

%================================================
% ACTO VI: EP / NTA
%================================================

\begin{frame}{Acto VI: Economía Poblacional y NTA (por qué esto es EP)}
Salud como transferencia en especie:
\[
\sum_x N(x)\tau_y(x) = \sum_x N(x)\tau_H(x)
\]

\bigskip
Transición demográfica y epidemiológica:
\[
G = \sum_x N(x)\Big[p_C(x)M_C(x)+p_N(x)M_N(x)\Big]
\]

\begin{itemize}
  \item envejecimiento: \(N(65+)\uparrow\)
  \item crónicas: \(p_N(x)\uparrow\), \(M_N(x)\uparrow\)
\end{itemize}

\bigskip
\emph{El diseño del financiamiento hoy redistribuye cargas y bienestar entre cohortes.}
\end{frame}

%------------------------------------------------
\begin{frame}{Cierre: la policy en una página}
\textbf{Problema:} cobertura amplia con financiamiento fragmentado + informalidad + poco espacio fiscal.

\bigskip
\textbf{Trade-off central:}
\[
\text{Equidad (}\beta\downarrow\text{)} \quad \text{vs} \quad \text{Incentivos (}\phi\uparrow\text{)}
\]

\bigskip
\textbf{Paquete de política creíble:}
\begin{enumerate}
  \item Paquete básico universal explícito (\(\alpha_0\)) para bajar \(\beta\).
  \item Diferencial contributivo claro (\(\Delta\alpha\)) para sostener \(\phi\).
  \item Privado como complemento regulado (evitar sistema dual).
  \item Sendero fiscal gradual y defendible (no “saltos”).
\end{enumerate}

\bigskip
\centering
\Large
Reforma de salud = diseño de contrato social bajo restricción fiscal.
\end{frame}

\end{document}